\documentclass[10pt,a4paper,sans]{moderncv}

\moderncvstyle{banking}
\moderncvcolor{blue}
\usepackage[scale=0.86,top=1.15cm,bottom=1.25cm]{geometry}
\usepackage{microtype}
\usepackage{enumitem}
\usepackage{multicol}
\usepackage{xcolor}
\usepackage{tikz}
\definecolor{deepblue}{RGB}{24,72,112}
\colorlet{color1}{deepblue}
\colorlet{color2}{black!62}
\setlength{\hintscolumnwidth}{2.55cm}
\setlist[itemize]{leftmargin=1.25em,itemsep=0.15em,topsep=0.2em,parsep=0pt,partopsep=0pt}
\renewcommand*{\sectionfont}{\Large\bfseries}
\newcommand*{\orcidicon}{%
  \tikz[baseline=-0.55ex]{%
    \fill[fill={rgb,255:red,166;green,206;blue,57}] (0,0) circle (0.72ex);%
    \node[white,font=\sffamily\bfseries,inner sep=0pt,scale=0.38] at (0,0) {iD};%
  }%
}
\AtBeginDocument{%
  \renewcommand*{\orcidsocialsymbol}{\orcidicon~}%
  \recomputelengths
}

\name{Rongfeng}{Cheng}
\title{Medical AI Research}
\address{Guangzhou, China}{}{}
\email{eb1287945452@mail.scut.edu.cn}
\homepage{crf2004.github.io/en}
\social[orcid]{0009-0002-3021-9647}

\newcommand{\hospitalaffiliation}{Institute of Healthcare Artificial Intelligence Application, Guangdong Second Provincial General Hospital}
\newcommand{\researchentry}[4]{%
  \noindent\begin{minipage}[t]{\textwidth}\vspace{0pt}
  \begin{minipage}[t]{0.78\textwidth}\vspace{0pt}\raggedright\textbf{#2}\end{minipage}\hfill
  \begin{minipage}[t]{0.20\textwidth}\vspace{0pt}\raggedleft\textit{#1}\end{minipage}\par
  \if\relax\detokenize{#3}\relax\else\vspace{0.12em}\textit{#3}\par\fi
  \vspace{0.12em}#4\par\vspace{0.48em}
  \end{minipage}\par
}

\begin{document}
\makecvtitle
\vspace{-1.2em}

\section{Education}
\cventry{Sep 2023--Jun 2027}{Bachelor of Management in Big Data Management and Application}{South China University of Technology (SCUT)}{Guangzhou, China}{}{Cumulative GPA: 3.46/4.0; average score: 85.18/100 (through Semester 6).\\Selected AI coursework: Large Models and Generative AI (93/100), Large Language Models and Prompt Engineering (93/100), Machine Learning (92/100), and Deep Learning (88/100).}

\section{Selected Achievements \& Deployed Work}
\cvitem{Published}{Co-authored a study on ethical safety in LLM code generation, published in \textbf{\textit{iScience}} (2026).}
\cvitem{IP}{Two first-inventor patent applications and two software copyrights.}
\cvitem{Clinical Deployment}{Co-developed Dingbei Xinxin, deployed at Guangdong Second Provincial General Hospital; led the five-person Dingbei Meimei engineering and clinical-workflow team through its hospital launch; led the data-layer development for a health-management product, including a standardized schema, field governance, and medical terminology mapping for multi-source health-examination data.}
\cvitem{Awards \& Projects}{Campus Gold, Internet+ Competition; Second Prize, Guangzhou Mathematical Modeling University League; Third Prize, CUMCM Guangdong Division; completed the SCUT SRP Student Research Training Project as Project Lead.}

\section{Research Experience}
\noindent\textit{\hospitalaffiliation}\par
\researchentry{Jun 2026--Present}{DLEF: Dataset-Label Evidence Fingerprints for Explaining Cross-Dataset Performance in ECG Classification}{First author; submitted to \textbf{IEEE BIBM 2026}}{\begin{itemize}
  \item \textbf{Problem:} Nominally identical diagnostic labels can rely on different ECG evidence across centers; standard performance metrics reveal external degradation without explaining its source.
  \item \textbf{Method:} Proposed and led DLEF, which applies controlled perturbations along rhythm, conduction, QRS morphology, lead-specific QRS, and repolarization axes to form dataset-label fingerprints and directional source-target distances.
  \item \textbf{Finding:} Across three public datasets and six diagnoses, DLEF distance closely tracked label-level external AUROC degradation (Spearman $r=0.843$) and outperformed prevalence, internal-performance, and embedding-shift baselines.
\end{itemize}}

\researchentry{Jun 2026--Present}{A Matched-Control Perturbation Audit Framework for Evaluating Evidence Use in Deep Learning Models: An Application to 12-Lead ECG Classification}{First author; revision in progress at \textbf{IEEE JBHI}}{\begin{itemize}
  \item \textbf{Problem:} Models can achieve strong discrimination while relying on evidence that conflicts with medical knowledge or becomes unstable under dataset shift.
  \item \textbf{Method:} Proposed and led an audit using domain-defined evidence units and matched controls, operator comparisons, and frozen-weight external transfer to test whether target effects exceed the strongest control and remain stable across datasets.
  \item \textbf{Finding:} Matched controls overturned the target-only 1AVB conclusion, operator choice changed AFIB and CRBBB verdicts, and external transfer exposed unstable evidence dependence; claims were classified as credible, limited, or unsupported.
\end{itemize}}

\newpage
\section{Internship \& Product Leadership}
\cventry{Nov 2024--Present}{LLM Development Engineer Intern}{AI Research Institute, Guangdong Second Provincial General Hospital}{}{}{\begin{itemize}
  \item Developed a standardized-patient Q\&A pipeline from approximately 8,000 cardiovascular case PDFs, covering LLM-based information extraction, entity disambiguation, Neo4j knowledge representation, and hybrid retrieval.
  \item Explored knowledge-graph paths as structural scaffolds for generating medical question--answer and reasoning-trace data for model fine-tuning, evaluation, and retrieval-augmented generation.
  \item Led full-stack and clinical-workflow development for the five-person Dingbei Meimei acne and facial-health assessment team; supported its initial release when the hospital's Refractory Acne Diagnosis and Treatment Center opened in July 2026.
\end{itemize}}

\cventry{Jul 2025--Jan 2026}{Backend \& Algorithm Development Intern}{Guangzhou Yunzhan Zhichuang Information Technology Co., Ltd.}{}{}{\begin{itemize}
  \item Designed a UMLS-grounded schema and normalization pipeline for heterogeneous health-examination data extracted from multi-institutional reports, including LLM-assisted schema evolution and terminology mapping.
  \item Developed backend workflow logic for the Health Navigation Plan within the company's iTwins product.
\end{itemize}}

\section{Products \& Team Projects}
\researchentry{2025--Present}{Dingbei Xinxin / CardioGPT: Cardiovascular Proactive Health LLM}{Co-founder \& Core Developer; Development Lead since Aug 2026}{\begin{itemize}
  \item Co-developed a cardiovascular health-management platform integrating PPG/ECG data, intelligent consultation, risk assessment, and proactive health workflows.
  \item Led model fine-tuning, RAG, and full-stack implementation using LoRA, DPO, MQ-RAG, and chain-of-thought methods; also prepared grant and competition application materials.
  \item Deployed the initial mini-program at Guangdong Second Provincial General Hospital and currently leads full-stack development of version 2.0; supported by national-level innovation-program projects in 2025 and 2026 and incorporated as Yujian Future Mind Technology (Zhongshan) Co., Ltd. in Aug 2026.
\end{itemize}}

\researchentry{2025}{Myocardial Ischemia Knowledge Graph}{Project Lead, SCUT Student Research Program}{Advisor: Xunde Dong. Led a five-student interdisciplinary team; designed the ontology, annotation schema, guidelines, and quality-control workflow; built a multi-annotator platform with adjudication and agreement checks; and implemented XGBoost-based implicit reasoning with graph-mapped interpretability analysis.}

\section{Publications \& Manuscripts}
\cvitem{Published [1]}{J. He, S. Qin, H. Zou, et al., \textbf{R. Cheng}, H. Zhang, and W. Cheng. ``CodeSafetyBench: Evaluating and Improving Ethical Safety in Large Language Model Code Generation.'' \textbf{\textit{iScience}}, 2026. \href{https://doi.org/10.1016/j.isci.2026.115073}{doi:10.1016/j.isci.2026.115073}.}
\cvitem{Revised [2]}{``Implicit Medical Ethical Deficits in Large Language Models: A Benchmarking and Alignment Study.'' \textbf{\textit{Nature Communications}}; led the revision-stage experiments for the response to reviewers.}
\cvitem{Under Review [3]}{\textbf{R. Cheng et al.} ``DLEF: Dataset-Label Evidence Fingerprints for Explaining Cross-Dataset Performance in ECG Classification.'' \textbf{\textit{IEEE BIBM 2026}}. First author.}
\cvitem{Revision in Progress [4]}{\textbf{R. Cheng et al.} ``A Matched-Control Perturbation Audit Framework for Evaluating Evidence Use in Deep Learning Models: An Application to 12-Lead ECG Classification.'' \textbf{\textit{IEEE JBHI}}. First author.}
\cvitem{Ongoing}{\textbf{PathCLIP: Evidence-Aware ECG--Text Representation Learning.} ECG--text alignment and zero-shot recognition with matched-count, permuted-hierarchy, residual-only, and held-out-label controls.}

% Keep the Skills taxonomy aligned with the web CV.
\section{Skills}
\begin{multicols}{2}
\cvitem{Programming \& Frameworks}{Python, PyTorch, scikit-learn, TypeScript, SQL}
\cvitem{Modeling Methods}{XGBoost, representation learning, multimodal learning, LLM fine-tuning, optimization}
\cvitem{LLM \& Knowledge Systems}{LoRA, DPO, RAG/MQ-RAG, knowledge graphs, information extraction, UMLS}
\columnbreak
\cvitem{Evaluation \& Data Curation}{RAG/retrieval evaluation, perturbation analysis, external validation, schema design}
\cvitem{Engineering}{Neo4j, Next.js, Flask, Git, Docker}
\end{multicols}
\vspace{-0.8em}

\section{Languages}
\cvitem{}{Chinese (native) \quad Cantonese (native) \quad English: IELTS 6.5, CET-6 528}

\end{document}
